%Across the seven epoch durations, both Fp1 and Fp2 maintained consistent pooled performance and similarly stable performance on Cao2018, although both electrodes' Internal results declined significantly relative to the 30~s reference at the longest durations (Figure~\ref{fig:f1_by_epoch}).

%Pooled macro $F_1$ varied by only 0.41 percentage points for Fp1, from 82.24\% at 120~s to 82.66\% at 20~s, and by 0.50 percentage points for Fp2, from 81.79\% at 120~s to 82.28\% at 20~s. None of the six alternative durations differed significantly from the 30~s reference after Bonferroni correction for either electrode, and the smallest corrected $p$-value was observed at 120~s for both (Fp1 $p=0.102$; Fp2 $p=0.098$).

Internal showed a less consistent, and electrode-dependent, pattern. The 30~s reference yielded each electrode's highest Internal $F_1$, 86.32\% for Fp1 and 86.82\% for Fp2. Relative to it, Fp1's $F_1$ fell significantly to 85.92\% at 50~s ($p=0.002$), to 85.67\% at 60~s ($p<0.001$) and to 85.04\% at 120~s ($p<0.001$). Fp2 differed significantly only at 120~s, falling to 85.57\% ($p<0.001$). Fp2's 50~s and 60~s declines, to 86.51\% and 86.27\% respectively, were descriptively similar to Fp1's but did not reach significance after correction, with $p=0.080$ and $p=0.064$ respectively. Neither electrode differed from its reference at 10, 20 or 40~s.
Significance tracked the consistency of the change rather than its magnitude.
%For Fp1, the 0.40-percentage-point mean drop at 50~s was significant because it recurred across sessions, with 35 of 46 falling below the reference and a median change of $-0.44$ percentage points. The larger 0.61-percentage-point mean drop at 10~s was not significant ($p=1.000$), because sessions split 22 below against 24 above the reference and the median session improved by 0.13 percentage points, so the mean was carried by a minority of sessions with large declines. The same pattern held for Fp2.

Performance on Cao2018 was stable for both electrodes, with no significant difference from the 30~s reference at any tested duration. For Fp1, the 30~s reference $F_1$ was 79.68\%, the lowest value across the seven durations. At the two shorter durations, the $F_1$ score were higher than the 30~s reference value, reaching 79.81\% at 10~s ($p=1.000$) and 79.90\% at 20~s ($p=0.381$). The $F_1$ is in its highest value, 80.05\% at 60~s, a peak 0.38 percentage points above the reference despite not reach significance ($p=0.882$). For Fp2, the 30~s reference $F_1$ was 78.45\%, also its lowest value. $F_1$ showed the same pattern, rising to 78.69\% at 10~s ($p=1.000$) and 78.71\% at 20~s ($p=0.159$), and reaching its highest value, 78.78\% at 120~s, a peak 0.33 percentage points above the reference that also did not reach significance ($p=1.000$).

%The evidence therefore supports stable Internal performance from 10 to 40~s for both electrodes, though not equivalence on Internal at 50~s and beyond for Fp1, or at 120~s for Fp2, and supports stable Cao2018 performance over the full 10--120~s range for both Fp1 and Fp2.
